\chapter{Algorithm Design, Computational Models, and Memory Hierarchy}
\label{cha:lecture-one}

\begin{center}
  \textbf{Lecture 1}\\
  \emph{Introduction, computational models, and memory hierarchy}\\[0.5em]
  Date: September 15, 2026
\end{center}

\section{Formal Definition of an Algorithm}
\label{sec:formal-algorithm}

An \emph{algorithm} is a finite and ordered sequence of unambiguous,
elementary instructions that, starting from a specified set of input data,
produces a corresponding set of output data within a finite amount of time.

In the classical formulation by Donald Knuth in \emph{The Art of Computer
Programming}~\cite{knuth1997}, an algorithm must satisfy five fundamental
properties:

\begin{description}
  \item[Finiteness.] The algorithm must always halt after a finite number of
    computation steps for every admissible input configuration.
  \item[Definiteness.] Each instruction must be stated clearly, precisely, and
    without ambiguity; the next execution state must be uniquely determined.
  \item[Input.] The algorithm receives zero or more quantities specified from
    outside before execution begins.
  \item[Output.] The algorithm produces one or more output quantities that
    have a specified relationship with the input data. The output must be
    meaningful for the problem being solved.
  \item[Effectiveness.] Every prescribed operation must be elementary and
    physically executable in finite time using a bounded amount of resources,
    such as paper and pencil or a sequential machine.
\end{description}

\subsection{Exact Algorithms, Artificial Intelligence, and Heuristics}
\label{sec:exact-vs-heuristic}

A classical exact algorithm is distinguished from an AI or heuristic method
by the guarantees it provides and by how it explores the space of decisions.

\begin{table}[htbp]
  \centering
  \caption{Exact algorithms compared with AI and heuristic methods.}
  \label{tab:algorithm-comparison}
  \begin{tabularx}{\textwidth}{@{}>{\raggedright\arraybackslash}p{0.20\textwidth}YY@{}}
    \toprule
    Dimension & Classical exact algorithm & AI or heuristic method \\
    \midrule
    Correctness guarantee & Produces a correct solution for every admissible
      instance; optimality is guaranteed when the problem specification asks
      for an optimum. & May prioritize empirical quality or practical
      efficiency and may lack a universal guarantee for every instance. \\
    Termination & Can be proved in advance using potential functions or loop
      variants. & Often tied to heuristic stopping criteria, such as a maximum
      number of epochs or gradient convergence. \\
    Decision space & Determined by a precisely specified procedure, which may
      be deterministic or randomized. & Explored through heuristic search,
      statistical inference, or loss-function optimization. \\
    \bottomrule
  \end{tabularx}
\end{table}

\emph{Heuristic programming} does not seek a formal proof of absolute
optimality over every admissible instance. Instead, it seeks an efficient
trade-off for computationally intractable problems, such as NP-complete or
NP-hard problems.

\section{Efficiency Measurement and Asymptotic Analysis}
\label{sec:asymptotic-analysis}

\subsection{Why Not Count the Exact Number of Instructions?}

Defining the exact running cost $T_A(n)$ as the actual number of clock cycles
or elementary instructions presents several difficulties:

\begin{enumerate}
  \item \textbf{Architecture dependence:} the result would depend on the
    specific CPU, instruction set architecture (ISA), and clock frequency.
  \item \textbf{Toolchain dependence:} it would vary with the compiler,
    optimization level, and operating system.
  \item \textbf{Lack of generality:} it would obscure the intrinsic quality of
    the algorithmic idea.
\end{enumerate}

For this reason, algorithm analysis uses asymptotic notation. It abstracts
away implementation details and multiplicative constants, and studies the
growth rate of the cost function as the input size tends to infinity:
$n \to \infty$.

\begin{description}
  \item[Big-$\mathcal{O}$ (asymptotic upper bound):]
    \[
      \begin{aligned}
        f(n) = \mathcal{O}(g(n)) &\iff
        \exists c > 0,\ n_0 > 0\text{ such that}\\
        &0 \leq f(n) \leq c\,g(n),\quad \forall n \geq n_0.
      \end{aligned}
    \]
  \item[Big-$\Omega$ (asymptotic lower bound):]
    \[
      \begin{aligned}
        f(n) = \Omega(g(n)) &\iff
        \exists c > 0,\ n_0 > 0\text{ such that}\\
        &0 \leq c\,g(n) \leq f(n),\quad \forall n \geq n_0.
      \end{aligned}
    \]
  \item[Big-$\Theta$ (tight asymptotic bound):]
    \[
      f(n) = \Theta(g(n)) \iff
      f(n) = \mathcal{O}(g(n)) \land f(n) = \Omega(g(n)).
    \]
\end{description}

\subsection{Why Prefer Worst-Case Analysis?}

\begin{itemize}
  \item \textbf{Absolute guarantee:} an upper bound guarantees that the
    algorithm will never take more time than the theoretical bound. This is
    essential for real-time and mission-critical systems.
  \item \textbf{Simple formal derivation:} worst-case analysis does not
    require probabilistic assumptions about the input distribution, which are
    often needed for average-case analysis.
\end{itemize}

\section{Computational Models: RAM and Real Hardware}
\label{sec:ram-model}

An algorithm is not executed in an abstract vacuum: it runs on a machine with
processors, memories, and data-transfer mechanisms. A computational model
specifies which machine details are relevant to the analysis and which are
deliberately ignored. The RAM model is the standard starting point; real
hardware then explains when that abstraction becomes insufficient.

\subsection{The Random Access Machine Model}

The Random Access Machine (RAM) model is a sequential, word-based abstraction
inspired by the von Neumann architecture. A RAM consists of a processor, a
finite program, and an addressable memory of words. Its usual assumptions are:

\begin{itemize}
  \item \textbf{Sequential execution:} one processor executes one elementary
    instruction at a time.
  \item \textbf{Word-sized data:} arithmetic, comparison, assignment, and
    address calculation on a constant number of words take constant time.
  \item \textbf{Uniform memory access:} reading or writing any address costs
    $\mathcal{O}(1)$, independently of the address and of the access history.
  \item \textbf{Sufficient memory:} the input and the data structures fit in
    the modeled address space; the cost of allocating the space is analyzed
    separately when it matters.
\end{itemize}

The RAM model is useful because it makes algorithmic ideas comparable. For
example, a loop that examines $n$ array entries performs $\Theta(n)$ RAM
operations, regardless of the particular processor used to run it. The model
does not claim that a real machine has infinite memory or truly constant
latency; it states which costs are being abstracted away.

\subsection{Real Hardware: Memory Hierarchy and the Memory Wall}

On a modern machine, the processor is much faster than the memories that hold
large datasets. This growing gap between computation and data movement is
known as the \emph{memory wall}. Moreover, local memory is not a single
uniform resource: it is organized into levels with different capacities,
latencies, and transfer granularities.

A simplified hierarchy, ordered from the fastest and smallest level to the
slowest and largest one, is:

\[
\begin{aligned}
  \text{CPU registers}
    &\longrightarrow \text{L1 cache}
    \longrightarrow \text{L2 cache}
    \longrightarrow \text{L3 cache}
    \longrightarrow \text{DRAM} \\
    &\longrightarrow \text{NVMe/SSD/disk}
    \longrightarrow \text{remote storage / cloud}.
\end{aligned}
\]

Registers and caches are part of the processor's fast memory system; DRAM is
the usual main memory; SSDs and disks provide persistent storage. Remote
storage and cloud services are included here only as a broader data-access
hierarchy, because they are reached through a network rather than through the
local memory bus.

\begin{figure}[htbp]
  \centering
  \begin{tikzpicture}[
    level/.style={draw, rounded corners=2pt, minimum height=8mm,
      align=center, font=\small},
    link/.style={-{Latex[length=2.5mm]}, thick},
    node distance=6mm
  ]
    \node[level, fill=blue!18, minimum width=28mm] (registers)
      {Registers};
    \node[level, fill=blue!14, minimum width=43mm, below=of registers] (l1)
      {L1 cache};
    \node[level, fill=blue!11, minimum width=58mm, below=of l1] (l2)
      {L2 cache};
    \node[level, fill=blue!8, minimum width=73mm, below=of l2] (l3)
      {L3 cache};
    \node[level, fill=green!12, minimum width=88mm, below=of l3] (dram)
      {DRAM / main memory};
    \node[level, fill=orange!16, minimum width=103mm, below=of dram] (storage)
      {NVMe / SSD / disk};
    \node[level, fill=gray!18, minimum width=118mm, below=of storage] (network)
      {Network / cloud};

    \draw[link] (registers) -- (l1);
    \draw[link] (l1) -- (l2);
    \draw[link] (l2) -- (l3);
    \draw[link] (l3) -- (dram);
    \draw[link] (dram) -- (storage);
    \draw[link] (storage) -- (network);

    \node[align=center, above=of registers, font=\footnotesize] (fast)
      {faster\\smaller\\more expensive per bit};
    \node[align=center, below=of network, font=\footnotesize] (large)
      {slower\\larger\\cheaper per bit};
  \end{tikzpicture}
  \caption{A hierarchical view of memory and storage.}
  \label{fig:memory-hierarchy}
\end{figure}

\autoref{fig:memory-hierarchy} describes the main local-memory levels:
each lower layer trades lower access speed for greater capacity.

Moving downward in \autoref{fig:memory-hierarchy}, capacity generally grows
while latency and transfer cost grow as well. The exact ratios depend on the
hardware and workload, so values such as $100$--$200$ times for DRAM relative
to L1, or $10^4$--$10^6$ times for disk relative to L1, should be read as
order-of-magnitude examples rather than universal constants.

\subsection{When the RAM Model Is Not Enough}

The RAM model is a good first approximation when the computation is mostly
sequential, the working set fits in main memory, and the cost of data movement
is not the dominant factor. A different model is preferable when one of the
following effects controls performance:

\begin{description}
  \item[GPU computing.] SIMT (Single Instruction, Multiple Threads) execution
    uses many processing cores and distinguishes shared, local, and global
    memory spaces.
  \item[External memory.] In out-of-core computation, the dataset does not
    fit in fast memory, so block transfers and I/O operations dominate the
    cost.
  \item[Parallel and distributed computation.] Several processors or machines
    execute simultaneously, and synchronization and communication become
    explicit costs.
  \item[Quantum computing.] Qubits, superposition, entanglement, and unitary
    gates require a model based on quantum operations rather than RAM words.
\end{description}

The engineering lesson is therefore not to discard the RAM model, but to use
the simplest model that still captures the bottleneck of the target system.

\section{Locality and the Two-Level Model}
\label{sec:external-memory}

Modern systems exploit two statistical properties of real programs:

\begin{enumerate}
  \item \textbf{Temporal locality:} if a memory location is referenced at time
    $t$, it is likely to be referenced again soon, at $t + \Delta t$.
  \item \textbf{Spatial locality:} if a memory location $x$ is referenced,
    nearby locations such as $x+1$, $x+2$, and so on are likely to be
    referenced soon.
\end{enumerate}

\subsection{The Two-Level Memory Model}

When processing gigabytes or terabytes of data, the multilevel hierarchy can
be abstracted as a two-level model, associated with the external-memory model
of Aggarwal and Vitter~\cite{aggarwal1988}, also known as the Disk Access
Machine model.

\begin{figure}[htbp]
  \centering
  \begin{tikzpicture}[
    box/.style={draw, minimum width=38mm, minimum height=10mm, align=center},
    arrow/.style={-{Latex[length=2.5mm]}, thick},
    node distance=12mm
  ]
    \node[box] (cpu) {CPU};
    \node[box, below=of cpu] (fast) {Fast memory $M$\\capacity: $M$ words};
    \node[box, below=of fast] (slow) {External memory $D$\\unbounded, blocks of $B$ words};
    \draw[arrow] (cpu) -- node[right, font=\small] {one word, $\mathcal{O}(1)$} (fast);
    \draw[arrow] (fast) -- node[right, font=\small] {one block = $B$ words} (slow);
  \end{tikzpicture}
  \caption{Two-level memory model with fast memory $M$ and external memory $D$.}
  \label{fig:two-level-memory}
\end{figure}

\autoref{fig:two-level-memory} describes the two-level abstraction: the
CPU operates on words in fast memory $M$, while transfers between $M$ and the
external memory $D$ occur in blocks of $B$ words.

\subsubsection{Model Parameters}

\begin{description}
  \item[$N$:] total number of problem elements to process.
  \item[$M$:] size of internal fast memory, measured as the maximum number of
    elements that can be held simultaneously.
  \item[$B$:] size of the block transferred in one I/O operation, with
    $M \geq 2B$.
\end{description}

\subsubsection{Cost Metrics}

\begin{enumerate}
  \item \textbf{I/O complexity (primary metric):} number of transfers of
    blocks of size $B$ between slow and fast memory.
  \item \textbf{CPU time (secondary metric):} number of logical and
    arithmetic operations performed after data have been loaded into $M$.
  \item \textbf{Space:} number of blocks allocated on secondary storage.
\end{enumerate}

This abstraction is fractal: it can model both RAM--disk transfers and
L3-cache--RAM transfers by redefining the meaning of $M$ and $B$.

\section{Worked Example: Linear Array Scan}
\label{sec:linear-scan}

Consider a contiguous array of $N$ integers scanned in sequence, for example
to compute the sum $\sum_{i=1}^{N} A[i]$.

\begin{lstlisting}[style=code, language=C, caption={Sequential array scan}, label={lst:linear-scan}]
long sum = 0;
for (int i = 0; i < N; i++) {
    sum += A[i];
}
\end{lstlisting}

\subsection{Analysis in the Classical RAM Model}

\begin{itemize}
  \item Memory accesses: $N$.
  \item Addition instructions: $N$.
  \item Asymptotic computation time: $\Theta(N)$.
\end{itemize}

\subsection{Analysis in the Two-Level Model}

Assume that the secondary memory stores array $A$ consecutively in blocks of
size $B$ and that $A[0]$ is aligned with a block boundary. Without this
alignment assumption, the scan may require at most one additional block
transfer, without changing the asymptotic bound:

\begin{enumerate}
  \item At the first access to $A[0]$, a miss occurs. The system transfers
    the block containing $A[0], \ldots, A[B-1]$ into fast memory with one I/O.
  \item The following $B-1$ accesses find their data already in fast memory,
    requiring no additional I/O. This exploits spatial locality.
  \item A new I/O transfer is needed only when the index crosses a block
    boundary, that is, once every $B$ elements.
\end{enumerate}

Therefore, the total I/O cost of a linear scan is
\[
  \operatorname{Scan}(N) = \left\lceil \frac{N}{B} \right\rceil
  = \Theta\!\left(\frac{N}{B}\right)\text{ I/O operations}.
\]

\subsection{Linear Access versus Sparse Access}

\begin{itemize}
  \item \textbf{Sequential access (high spatial locality):} reading $N$
    elements requires $N/B$ I/O operations. Since $B$ typically ranges from
    $64$ bytes for a cache line to hundreds of kilobytes or megabytes for
    pages and disk blocks, the speedup over the naive model can be a factor of
    approximately $B$.
  \item \textbf{Random access or stride $k>B$ (no spatial locality):} every
    requested element may reside in a different block, causing $N$ independent
    I/O operations. The cost degrades to $\Theta(N)$ I/Os, which is
    asymptotically slower by a factor of $B$.
\end{itemize}
